\documentclass[11pt]{article}
\usepackage[margin=25mm]{geometry}
\usepackage{fontspec}
\setmainfont{Latin Modern Roman}
\usepackage{amsmath,amssymb,mathtools}
\usepackage{unicode-math}
\setmathfont{Latin Modern Math}
\usepackage{xurl,hyperref}
\hypersetup{colorlinks=true,urlcolor=blue}
\setcounter{secnumdepth}{0}
\setlength{\parindent}{0pt}
\setlength{\parskip}{5pt}
\setlength{\emergencystretch}{3em}
\newcommand{\sref}[2]{\hyperlink{source:#1:#2}{[S#2]}}
\newcommand{\cataloguescope}[1]{\paragraph{Recorded target.}#1}
\newcommand{\R}{\mathbb R}
\newcommand{\C}{\mathbb C}
\newcommand{\N}{\mathbb N}
\newcommand{\Z}{\mathbb Z}
\newcommand{\Q}{\mathbb Q}
\newcommand{\F}{\mathbb F}
\newcommand{\eps}{\varepsilon}
\begin{document}
\section{Renormalization through Langmann–Szabo Duality}\label{statement}

\subsection{Definitions and mathematical statement}

\paragraph{Shared notation.}$\mathbb N=\{1,2,\ldots\}$; the standard number fields have their usual meanings. All additional conventions are stated locally.


\paragraph{Definitions and conventions.} The Moyal deformation replaces ordinary multiplication of fields by a noncommutative star product with coordinate commutators $[x^\mu,x^\nu]=i\theta^{\mu\nu}$. Langmann--Szabo duality relates position and momentum descriptions of certain such scalar models, exchanging the kinetic and harmonic-oscillator terms. Perturbative renormalizability means that ultraviolet divergences at all perturbative orders can be absorbed into a fixed finite set of parameters and field normalizations.

\paragraph{Source scope and qualifications.} The original report already describes the successful harmonic scalar construction. Its broad construction question is historical and model dependent; it must not be converted into a claim that no example is known, or into an unsupported theorem for all noncommutative field theories.

\paragraph{Catalogue mathematical target.}
\paragraph{Catalogue wording retained for traceability.}
Can Langmann–Szabo duality be used to construct renormalizable noncommutative quantum field theories?

\paragraph{Saved literature context.} The supplied catalogue's assessment dated 2026-08-21 records: Langmann--Szabo duality led to renormalizable harmonic noncommutative scalar theories, while the general construction program and gauge-theory extension remain model dependent. The assessment is retained as dated provenance; it is not a new comprehensive literature or open-status certification.
\subsection{Short English statement}
Use position--momentum duality to construct and understand renormalizable noncommutative quantum field models.
\subsection{Research attempt}
\paragraph{Model, date and investigated scope.}
OpenAI GPT-6 Astra (configured model \texttt{gpt-6-astra}), ultra reasoning; generated 11 October 2026. Derives the quadratic position-momentum exchange and shows why symmetry alone cannot imply finite-parameter renormalizability.

\paragraph{Sources and literature check.}
The complete supplied definition, including source qualifications and the dated literature assessment, was reviewed. The calculation below is an elementary restricted attempt motivated by [S2]; it does not claim novelty or a fresh literature-status determination.

\paragraph{Context and investigated scope.} The complete supplied definition, including source qualifications and the dated literature assessment, was reviewed. The calculation below is an elementary restricted attempt motivated by [S2]; it does not claim novelty or a fresh literature-status determination.

\paragraph{Quadratic duality test.} Fix $\Omega>0$. For a Schwartz function $f$ on $\mathbb R^d$, use the unitary Fourier transform $\widehat f(p)=(2\pi)^{-d/2}\int e^{-ipx}f(x)dx$. Integration by parts gives $\widehat{\partial_jf}=ip_j\widehat f$ and $\widehat{x_jf}=i\partial_{p_j}\widehat f$. Consequently
\[
 Q_\Omega(f)=\int(|\nabla f|^2+\Omega^2|x f|^2+m^2|f|^2),
 \quad Q_\Omega(\widehat f)=\Omega^2Q_{1/\Omega}(f)
\]
provided the mass coefficient on the right is replaced by $m^2/\Omega^2$; for $\Omega=1$ no replacement is needed. This is the oscillator-level mechanism behind exchanging position and momentum, in units where the exchange has unit scale.

\paragraph{Attempt to infer renormalizability.} The calculation selects a self-dual quadratic form, but self-duality does not restrict an action to finitely many interactions. For each integer $k\geq2$, the functional
\[
 I_k(f)=\int |f(x)|^{2k}dx+\int|\widehat f(p)|^{2k}dp
\]
is Fourier invariant, because two Fourier transforms reflect the argument. These functionals are independent: replacing $f$ by $tf$ gives distinct powers $t^{2k}$. Thus symmetry alone allows arbitrarily many couplings even on this simple scalar function space.

\paragraph{Boundary and source scope.} These $I_k$ are illustrative nonlocal functionals, not asserted to be admissible Moyal-local interactions. Their role is to invalidate the inference from duality alone to a finite counterterm space. The missing step is a model-specific all-order power-counting and subtraction argument respecting the star product. The report already contains successful harmonic scalar constructions; this exercise neither rediscovers their proof nor settles gauge extensions or all noncommutative theories.

\paragraph{Verification and adversarial audit.}
Checked Fourier normalization, the mass rescaling, the self-dual point, and independence of invariant functionals; distinguished nonlocal test interactions from the source model. Model self-audit only; no independent human proof checking. Prompts A and B were applied to this canonical record and its qualified source scope. Checked Fourier normalization, the mass rescaling, the self-dual point, and independence of invariant functionals; distinguished nonlocal test interactions from the source model. Retained unresolved status and zero original-target progress; A model-specific all-order counterterm analysis, especially for broader interactions or gauge models, is missing.

Coordinator checked Fourier derivative identities, mass rescaling and independence of the invariant functionals under amplitude scaling. Omega must be positive; this correction was sent to the author. This additional check was model-assisted, not independent human verification.

\paragraph{Tools, process and contributions.}
Python 3 was used for file preparation and structural checks; no numerical experiment or theorem prover was used. Contribution prompt A was followed by prompt B. GPT-6 Astra with ultra reasoning drafted the research and performed its stated self-audit; any additional coordinator contribution is identified in the verification and audit. A GPT-6-based Codex coordinator assembled the manuscript and checked repository integration; its exact variant and reasoning configuration were not exposed. No human mathematical contribution, independent human proof check or formal proof certification is supplied. Complete derivations and the consulted source references are included here; private reasoning transcripts are not published.

\paragraph{Final status and first remaining gap.}
\textbf{UNRESOLVED}. A model-specific all-order counterterm analysis, especially for broader interactions or gauge models, is missing.

\section{Completion Estimate}
\noindent\textbf{COMPLETION: 0\%}
The calculation is an auxiliary or restricted attempt and establishes no newly verified portion of the full source target. A model-specific all-order counterterm analysis, especially for broader interactions or gauge models, is missing.

\subsection{Sources}
\begin{itemize}
\item[\hypertarget{source:statement:1}{S1}] ULAM.AI and the UnsolvedMath Contributors, UnsolvedMath, catalogue entry. Catalogue curation under CC BY 4.0; definitions and English exposition added locally.
\url{https://huggingface.co/datasets/ulamai/UnsolvedMath}.
\item[\hypertarget{source:statement:2}{S2}] Noncommutative Geometry and Quantum Field Theory (2005), Oberwolfach Reports, DOI 10.4171/owr/2005/48, pp. 2705--2760; official report source and mathematical context.
\url{https://doi.org/10.4171/owr/2005/48}.
\item[\hypertarget{source:statement:3}{S3}] Official EMS Press PDF of the same Oberwolfach report; PDF pages 17--19.
\url{https://ems.press/content/serial-article-files/46022}.
\item[\hypertarget{source:statement:4}{S4}] Langmann--Szabo duality; reference in the saved dated catalogue literature review.
\url{https://arxiv.org/abs/hep-th/0202039}.
\item[\hypertarget{source:statement:5}{S5}] Renormalizable UV--IR duality models; reference in the saved dated catalogue literature review.
\url{https://arxiv.org/abs/1001.3776}.
\end{itemize}
\label{end:statement}
\end{document}
